\clearpage
\SourceLeaf{311}
\GaussVSourceUnnumbered
\GaussVSourceGathering{39}
\pdfbookmark[1]{Göttingische gelehrte Anzeigen. 1840 März 26.}{anzeigen-1840-maerz-26}
\begin{center}
\rule{.70\linewidth}{.4pt}\\[.7\baselineskip]
Göttingische gelehrte Anzeigen. 1840 März 26.\\[.7\baselineskip]
\rule{.70\linewidth}{.4pt}
\end{center}

Der Königl. Societät der Wissenschaften ist am 9. März von dem Hofrath
\textsc{Gauss} eine Vorlesung überreicht mit der Überschrift:
\begin{center}
\textit{Allgemeine Lehrsätze in Beziehung auf die im verkehrten Verhältnisse des Quadrats
der Entfernung wirkenden Anziehungs- und Abstossungskräfte,}
\end{center}
\noindent
deren Gegenstand hier, so weit die Bestimmung dieser Blätter es verstattet, näher
bezeichnet werden soll.

Wenn man zur Erklärung der magnetischen Erscheinungen zwei magnetische
Flüssigkeiten annimmt, deren gleichnamige Theile einander abstossen, die
ungleichnamigen einander anziehen, so besteht das Magnetisirtsein eines Körpers
in der \textit{Scheidung} der in ihm enthaltenen Flüssigkeiten. Nach dem ersten Eindrucke
des sinnlichen Scheines ist man geneigt, diese Scheidung an einem Magnetstabe
sich so vorzustellen, dass das eine Ende nur die nördliche Flüssigkeit, das andere
die südliche enthalte: allein genauere Überlegung zeigt bald die Unstatthaftigkeit
einer solchen Vorstellungsart, und die Nothwendigkeit, die Scheidung nur in den
kleinsten für uns unmessbaren Theilen des Trägers der Flüssigkeiten (Stahls oder
Eisens) anzunehmen, so dass jeder \textit{messbare} Theil des Trägers nach der Scheidung
wie vor derselben immer gleiche Quantität beider Flüssigkeiten enthält. Wenn
nun aber in physicalischen Schriften auf dem einen
\clearpage
\SourceLeaf{312}
\GaussVSourceFolio{306}{L}{ANZEIGEN.}
\noindent
Blatte die richtige Vorstellungsart gelehrt, und doch auf dem folgenden wieder
von dem freien nördlichen Magnetismus, der sich in dem einen Ende des Magnetstabes,
und dem südlichen, der sich in dem andern befinden soll, geredet wird,
so scheint eine solche schwankende Sprache die Begriffe zu verwirren, und
wissenschaftlicher Präcision abzusagen. In diese Unklarheit kann nur Licht
gebracht werden durch ein Theorem, welches in der \textit{Intensitas vis magneticae
terrestris} Art. 2 angekündigt ist, und darin besteht, dass anstatt irgend welcher
Vertheilung der magnetischen Flüssigkeiten innerhalb eines begrenzten körperlichen
Raumes substituirt werden kann eine ideale Vertheilung auf der Oberfläche
dieses Raumes mit dem Erfolge, dass von dieser idealen Vertheilung in jedem
Punkte des äussern Raumes genau dieselbe Wirkung ausgeübt wird, wie von jener
wirklichen. Auch in der \textit{Allgemeinen Theorie des Erdmagnetismus} (im dritten
Jahrgange der von dem Verf. gemeinschaftlich mit Hrn. Prof. \textsc{Weber} herausgegebenen
\textit{Resultate des magnetischen Vereins}) hat der Verf. sich auf dieses
Theorem bezogen, indem er bemerkt hat, dass es zwar unmöglich ist, die wirkliche
Vertheilung des Magnetismus im Innern der Erde zu erforschen, dass aber
die äquivalirende ideale Vertheilung auf der Erdoberfläche in unserm Bereiche
liegt; eine graphische Darstellung dieser idealen Vertheilung, nach der erwähnten
Theorie, ist bereits gezeichnet und lithographirt, und wird in dem magnetischen
Atlas mit enthalten sein, der in Kurzem mit dem vierten Theile der \textit{Resultate}
erscheinen wird.

Was nun die Begründung des in Rede stehenden Theorems betrifft, so
erfordert dieselbe eine ziemlich zusammengesetzte mathematische Zurüstung; das
Theorem selbst erscheint als ein specieller Fall eines allgemeinern, welches
seiner Seits das letzte Glied einer Kette genau zusammenhängender allgemeiner
Lehrsätze bildet. Die vollständige Entwickelung dieser Untersuchungen ist der
Gegenstand der vorliegenden Abhandlung. Es ist jedoch nicht die Meinung,
dass die Zwischensätze bloss wie Mittel zu dem angeführten Zwecke betrachtet
werden sollen, sondern sie nehmen als allgemeine Untersuchungen über die im
verkehrten Verhältnisse des Quadrats der Entfernung wirkenden Kräfte (wovon
die magnetischen nur ein einzelner Fall sind) ein selbstständiges Interesse in
Anspruch: allein in das Einzelne hier näher einzugehen, würde eine grössere
Ausführlichkeit erfordern, als der Raum dieser Blätter verstattet. Nur ein paar
der Lehrsätze, die ohne grosse Vorbereitung verständlich zu machen sind, und mit
\clearpage
\SourceLeaf{313}
\GaussVSourceFolio{307}{R}{ALLGEMEINE LEHRSÄTZE IN BEZIEHUNG AUF DIE IM VERKEHRTEN VERHÄLTNISSE ETC.}
\GaussVSourceGathering{39*}
\noindent
Untersuchungen anderer Mathematiker in einiger Berührung stehen, mögen hier
als Proben erwähnt werden.

I. Eine Gleichgewichtsfläche in Beziehung auf Massen, die anziehende oder
abstossende Kräfte ausüben, heisst bekanntlich jede Fläche, in deren sämmtlichen
Punkten die Resultante der Kräfte entweder gegen die Fläche normal ist, oder
selbst verschwindet. Eins der Theoreme ist nun folgendes: Wenn eine geschlossene
Fläche eine Gleichgewichtsfläche für die Anziehungs- oder Abstossungskräfte
von Massen ist, die sich sämmtlich im äusseren Raume befinden, so ist
die Resultante der Kräfte so wohl in jedem Punkte jener Fläche, als auch in
jedem Punkte des ganzen innern Raumes nothwendig \(=0\).

\textsc{Poisson} bemerkt in seiner berühmten Abhandlung über die Vertheilung der
Electricität an der Oberfläche leitender Körper, dass es zur Erhaltung eines
beharrlichen electrischen Zustandes eines electrisirten leitenden Körpers nicht
zureichend sei, dass die innere Grenzfläche der freien an der Oberfläche des
Leiters befindlichen Electricität eine Gleichgewichtsfläche sei, sondern noch
ausserdem erforderlich, dass diese Electricität auch in keinem Punkte des innern
Raumes Anziehung oder Abstossung ausübe.

Das oben erwähnte Theorem beweist dagegen, dass allerdings die erste
Bedingung allein hinreicht, in sofern sie die zweite als eine nothwendige Folge
schon in sich begreift.

II. Ein zweites Theorem bezieht sich auf den andern Fall, wo die anziehenden
oder abstossenden Massen sich innerhalb des von einer geschlossenen
Fläche begrenzten Raumes befinden. Hier wird in jedem Punkte der Fläche,
wenn sie eine Gleichgewichtsfläche ist, die resultirende Kraft nach Einerlei Seite
gerichtet sein, auch wenn anziehende und abstossende Massen zugleich vorhanden
sind; je nachdem nemlich das Aggregat der ersteren, oder das der anderen das
grössere ist, wird die Resultante in allen Punkten nach innen oder nach aussen
gerichtet sein: ist aber das Aggregat der anziehenden Massen dem der abstossenden
gleich, so wird, wenn es überhaupt eine geschlossene und einschliessende
Gleichgewichtsfläche gibt, die Resultante der Kräfte in jedem Punkte derselben,
und zugleich im ganzen unendlichen äussern Raume, \(=0\) sein.

III. In der Abhandlung ist ein strenger Beweis geführt, nicht bloss dafür,
dass auf jeder gegebenen geschlossenen Fläche eine gegebene Gesammtmasse so
nach der Stetigkeit vertheilt gedacht werden kann, dass in jedem
\clearpage
\SourceLeaf{314}
\GaussVSourceFolio{308}{L}{ANZEIGEN.}
\noindent
Punkte des innern Raumes die Resultante der Anziehungs- oder Abstossungskräfte
\(=0\) wird, sondern auch, dass dies allemal \textit{nur auf eine einzige Art} möglich ist.
Gerade das Gegentheil dieses Theorems war unlängst von einem geschickten
Geometer behauptet, in einer der Pariser Academie der Wissenschaften gemachten
Mittheilung (\textit{Comptes rendus} 1839. Nr. 6).
\begin{center}\rule{.26\linewidth}{.4pt}\end{center}
